\documentclass[11pt]{article}
\usepackage[margin=1in]{geometry}
\usepackage{amsmath,booktabs,graphicx,hyperref}
\title{A Verified Flat-Bed Comparison of Linear and Nonlinear Shallow-Water Models}
\author{JIajun FU}
\date{Project working draft, V0.2, October 2, 2026}
\begin{document}
\maketitle
\begin{abstract}
We implement and verify constant-depth linear and nonlinear shallow-water
solvers and compare their predictions under matched physical initial conditions.
Independent smooth exact solutions demonstrate approximately second-order mean
$L^1$ convergence. An amplitude sweep separates waveform model differences from
empirical mesh sensitivity. These results establish a computational baseline;
they do not constitute validation against observed tsunamis.
\end{abstract}
\section{Models and numerical methods}
With static depth $D$, free-surface displacement $\eta$, depth-averaged velocity
$u$, and $h=D+\eta$, the nonlinear conservation laws are
\[
h_t+(hu)_x=0,\qquad (hu)_t+(hu^2+gh^2/2)_x=0.
\]
Linearization about $h=D,u=0$ yields $\eta_t+Du_x=0$, $u_t+g\eta_x=0$.
Scaling $\eta$ by $A$, $u$ by $(A/D)\sqrt{gD}$ and time by
$L/\sqrt{gD}$ places a factor $\epsilon=A/D$ before the nonlinear terms.
Small $\epsilon$ does not independently ensure negligible accumulated error
or the validity of the nondispersive shallow-water approximation.

For $W=(\eta,u)^T$ and $B=\left(\begin{smallmatrix}0&D\\g&0\end{smallmatrix}\right)$,
the separately implemented exact upwind flux is
$\widehat F= B(W_l+W_r)/2-c(W_r-W_l)/2$, $c=\sqrt{gD}$.
The nonlinear solver uses HLL. Both flat-bed solvers use componentwise
MC-limited MUSCL reconstruction and SSP-RK2 at CFL 0.4. No evolved water-depth
clipping is applied. The variable-bed hydrostatic-reconstruction solver from
V0.1 remains first order and is excluded from this comparison.

\section{Verification and experimental controls}
Linear traveling waves and a smooth nonlinear simple wave provide independent
exact references, evaluated as cell averages with eight-point Gaussian quadrature.
For $D=10$ m, $A=0.2$ m, $L=1000$ m and $t=20$ s, refinement from 400 to
800 cells gives mean $L^1$ orders 1.9959 and 1.9985, respectively. Limited extrema
converge more slowly in the maximum norm. Nineteen automatic tests pass,
including conservation, boundary accounting and input-contract checks.

The amplitude sweep uses a periodic domain with
$\eta_0=A\sin(2\pi x/L)$ and $u_0=c\eta_0/D$, at $t=0.2L/c$.
Each model receives cell averages of its own evolved variables; nonlinear
momentum is averaged from $(D+\eta_0)u_0$, not $Du_0$.
This is a matched-continuous-field comparison; the nonlinear initial field is
not an exact nonlinear simple wave. Four nested meshes (400--3200 cells) are
used for each amplitude, with 6400 cells additionally used for $\epsilon=10^{-4}$.

\section{Results}
Define $E=\operatorname{RMS}(\eta_{NL}-\eta_L)/A$. Empirical mesh sensitivity
$S$ is the sum of each model's fine-to-coarse RMS differences divided by $A$,
using conservative restriction. The prespecified screen $E/S\geq5$ is not
a rigorous numerical error bound.
\begin{center}
\begin{tabular}{rrrr}
\toprule
$A/D$ & Finest cells & $100E$ & $E/S$\\
\midrule
0.0001 & 6400 & 0.00625148 & 12.78\\
0.001 & 3200 & 0.0625149 & 40.24\\
0.01 & 3200 & 0.625149 & 515.50\\
0.05 & 3200 & 3.12577 & 3342.28\\
0.1 & 3200 & 6.25171 & 6648.85\\
0.2 & 3200 & 12.5053 & 12714.27\\
\bottomrule
\end{tabular}
\end{center}
All waveform differences pass the empirical screen. The smallest amplitude
fails it at 3200 cells ($E/S\approx4$) and passes after refinement. For this
initial shape and time, $E\approx0.625\epsilon$. This coefficient is not
universal. The maximum normalized mass residual in the comparison is
$3.64\times10^{-16}$, which alone does not bound waveform error.
Crest-position differences at the two smallest amplitudes remain unresolved
by the same mesh screen and are not interpreted as resolved phase shifts.

\section{Limitations and next stage}
Both models omit dispersion and use a fully wet periodic flat bed. No coastal
shoaling, inundation, physical dissipation calibration or observational validation
is included. A subsequent study should first verify a well-balanced higher-order
variable-bed method and a variable-depth linear solver, then investigate
controlled shoaling and reflection before comparison with laboratory data.
This is an intermediate project report, not a claim of a novel numerical method.

\end{document}
